\clearpage
\begin{trapbox}
\textbf{误区一：把割线的对称，误当成交点的对称}

\begin{center}
\begin{tikzpicture}[x=1cm,y=1cm,font=\small]
\draw[gray!45,rounded corners=3pt,fill=white] (0,0) rectangle (13.8,8.7);
\node[anchor=west,font=\small\bfseries] at (.4,8.2) {割线关于 x=2 对称，不代表第二交点 M 与 N 对称};
\begin{scope}[shift={(1.25,4.15)},scale=1.25]\CNBase

\coordinate (S) at ({18/7},{-1/7});
\draw[purple!75!black,dashed] (M)--(S);
\draw[purple!75!black] ($(S)+(-.07,-.07)$)--($(S)+(.07,.07)$) ($(S)+(-.07,.07)$)--($(S)+(.07,-.07)$);
\node[purple!75!black,anchor=west] at (2.8,.2) {$M^*$};
\end{scope}
\node[anchor=west] at (8.1,7) {$M^*$ 才是 $M$ 的镜像};
\node[anchor=west,purple!75!black] at (8.1,5.65) {$M^*=(18/7,-1/7)$};
\node[anchor=west] at (8.1,4.3) {$M^*$ 在直线 $PN$ 上};
\node[anchor=west] at (8.1,3.05) {却不在双曲线上};
\node[anchor=west] at (8.1,1.65) {$N\ne M^*$，交点不对称};

\end{tikzpicture}

\end{center}

\textbf{错法：} 两条割线关于 $x=2$ 对称，就把 $M$ 的镜像当成 $N$。

\textbf{为什么错、怎样改：} 真正镜像点是 $M^*=(18/7,-1/7)$，它在 $PN$ 上，却不满足双曲线方程。$N=(26/7,-17/7)$ 是另一个对象。必须由“仍在曲线上”筛选交点，不能只用直线的反射定位。

\medskip
\textbf{误区二：平行就写成同一条线。} 本例弦是 $x+y=9/7$，切线是 $x+y=1$。应先设 $x+y=t$，再由其他条件求位置。

\medskip
\textbf{误区三：忘掉公共点。} 两条任意位置的直线即使斜率相反，也不满足结论；必须同过双曲线上的 $P$，各自另有不同于 $P$ 的交点。
\end{trapbox}

\clearpage
\begin{trapbox}
\textbf{误区四：把两种缺少第二交点的情况混为一谈}

\begin{center}
\begin{tikzpicture}[x=1cm,y=1cm,font=\small]
\draw[gray!45,rounded corners=3pt,fill=white] (0,0) rectangle (13.8,8);
\node[anchor=west,font=\small\bfseries] at (.4,7.5) {两种不同的失效原因：相切，与平行渐近线};
\begin{scope}[shift={(1,3.3)},scale=.9]
\begin{scope}\clip (0,-2.1) rectangle (5,3);
\draw[cncurve,samples=120,domain=-2:2.9,variable=\t] plot ({sqrt(2*(1+\t*\t))},\t);
\draw[cnsec] (0,-1)--(5,4);
\end{scope}
\fill (2,1) circle(2pt) node[below right=7pt] {$P$};
\end{scope}
\node at (3.2,6.9) {$m=1$};
\node[align=center] at (3.2,.85) {切线：重根仍是 $P$\\没有不同于 $P$ 的交点};
\begin{scope}[shift={(7.7,3.3)},scale=.9]
\begin{scope}\clip (0,-2.1) rectangle (5,3);
\draw[cncurve,samples=120,domain=-2:2.9,variable=\t] plot ({sqrt(2*(1+\t*\t))},\t);
\draw[cnsec] (0,-0.41421356237309515)--(5,3.121320343559643);
\draw[gray,dashed] (0,0)--(5,{5/sqrt(2)});
\end{scope}
\fill (2,1) circle(2pt) node[below right=7pt] {$P$};
\end{scope}
\node at (9.9,6.9) {$m=1/\sqrt2$};
\node[align=center] at (9.9,.85) {平行渐近线\\二次项消失，仅剩 $P$};

\end{tikzpicture}

\end{center}

\textbf{错法：} 在 $P$ 相切时，把重根当成新交点；或让割线平行渐近线，仍套第二交点公式。

\textbf{怎样分别检查：} 设增量 $h$，交点方程为
\[
h\left[Dh+2\left(\frac{u}{a^2}-\frac{vk}{b^2}\right)\right]=0.
\]
切线情形是括号常数项为零，只有重合的 $h=0$；平行渐近线是 $D=0$，方程没有所需的第二个有限根。二者原因不同，都不能当成题设割线。

\medskip
\textbf{误区五：顶点处除以零。} $v=0$ 时不能套 $-b^2u/(a^2v)$。例3的图已经给出替代判断：弦与切线同为竖直线，平行结论仍成立。
\end{trapbox}

\clearpage
\begin{trapbox}
\textbf{误区六：漏掉角平分线中的“内角”}

\begin{center}
\begin{tikzpicture}[x=1cm,y=1cm,font=\small]
\draw[gray!45,rounded corners=3pt,fill=white] (0,0) rectangle (13.8,8.2);
\node[anchor=west,font=\small\bfseries] at (.4,7.7) {斜率相反，还要检查射线朝向，才能判断内角平分线};
\foreach \xx/\kind in {3.5/1,10.3/2}{
\begin{scope}[shift={(\xx,4.5)}]
\draw[cnmirror] (0,-2)--(0,2.1);
\ifnum\kind=1
 \fill[green!45!black,opacity=.08] (0,0)--(-1,-2)--(1,-2)--cycle;
 \draw[cnsec,-Stealth] (0,0)--(-1,-2);
 \draw[cnsec,-Stealth] (0,0)--(1,-2);
 \node[left] at (-1,-2) {$PM$}; \node[right] at (1,-2) {$PN$};
 \node[align=center] at (0,-3) {$m=2$，射线分居轴两侧\\竖直轴平分内角};
\else
 \fill[orange!85!black,opacity=.08] (0,0)--(-2,-1)--(-2,1)--cycle;
 \draw[cnsec,-Stealth] (0,0)--(-2,-1);
 \draw[cnsec,-Stealth] (0,0)--(-2,1);
 \node[left] at (-2,-1) {$PM$}; \node[left] at (-2,1) {$PN$};
 \draw[gray] (0,0)--(-2.3,0);
 \node[align=center] at (0,-3) {$m=1/2$，射线位于轴同侧\\竖直轴平分外角};
\fi
\fill (0,0) circle(1.8pt); \node at (.5,.35) {$P$};
\end{scope}}
\node at (6.9,.5) {仅画例中两条边射线的方向，不表示交点到 $P$ 的距离。};
\end{tikzpicture}

\end{center}

\textbf{错法：} 由两条直线斜率相反，直接认定内心就在 $x=u$ 上。

\textbf{为什么错：} 左图取首页的 $m=2$，两个第二交点位于竖直轴两侧；右图取 $m=1/2$，实际第二交点为 $M(-2,-1),N(-10,7)$，都在 $x=2$ 左侧。两种情况下割线都对称，但竖直轴在右图平分的是外角。

\textbf{怎样改：} 从内心条件出发时，确认内心在三角形内部；反向判断时，继续检查边射线是否分居竖直轴两侧。只求弦方向的 C095 不要求该轴是内角平分线；内心例题则必须加上这一条件。
\end{trapbox}
